\textbf{Methods.} (i) $s\le5.24$: the spectral branch of \texttt{work/kkr-extremal}.
(ii) $2\le s\le8.6$: two-sided shooting (\texttt{strong/twoside.py}): outward from the horizon,
with its formal series (rank-one level blocks) plus the non-analytic throat mode $B\,x^\gamma$
($\gamma$ the exact throat exponent), inward from the large-$r$ series with free mass $\mu$ and
squashing amplitude $b_4$; $m=3\mu/4$. (iii) $s\ge9$: multiple shooting (\texttt{strong/ms.py},
\texttt{ms\_cont.py}). Single shooting fails for $s\gtrsim10$ because in the core the squashing
mode behaves as $e^{\pm4\sqrt{6}\alpha/r}$ (a factor $e^{76}$ at $\alpha=10$). The solver places
nodes at equal WKB phase in the core, integrates segments with $x\le0.5\rho$ in quadruple
precision ($\_\_$float128 Bulirsch--Stoer) from 60-digit horizon series, and continues in $\alpha$
across the resonances $\gamma=3,4$ ($\alpha=4.83,13.97$).

\textbf{Checks.} At the common point $s=2.98202$ the two-sided and spectral values differ by
3.4e-07; against a spline of the spectral branch the shooting/BVP
values deviate by at most 2.0e-05, the size of the spectral branch's own constraint
violations (median 3e-06, max 7e-05). Varying the shooting parameters moves $m$ by
at most 2.8e-08.
Multiple-shooting Newton residuals are below 3e-10. A degree-6 polynomial through the
18 multiple-shooting points with $9.16\le s\le10.3$, extrapolated to the BVP point $s=9.0037$,
agrees with it to 2e-09.

\begin{center}\small\begin{tabular}{rrrl}\toprule $s$ & $m(s)$ & $(m-3s/4)\sqrt s$ & method\\\midrule
2 & 2.4446986200 & 1.33600560 & shoot\\
2.4 & 2.6802239675 & 1.36363711 & shoot\\
2.98202 & 3.0410066937 & 1.38923766 & shoot\\
3.9 & 3.6394756194 & 1.41097629 & shoot\\
5.07 & 4.4346318566 & 1.42334984 & shoot\\
6.591 & 5.4998298751 & 1.42890434 & shoot\\
8.5683 & 6.9145757941 & 1.42948377 & shoot\\
9.00366 & 7.2290351282 & 1.42916958 & BVP\\
9.16046 & 7.3425005931 & 1.42903362 & MS\\
9.32 & 7.4580467159 & 1.42888461 & MS\\
9.52 & 7.6030391456 & 1.42868400 & MS\\
9.59986 & 7.6609741162 & 1.42859997 & MS\\
9.86 & 7.8498672728 & 1.42831221 & MS\\
10.08 & 8.0097946897 & 1.42805387 & MS\\
11.0299 & 8.7020175240 & 1.42681967 & MS\\
13.5121 & 10.5212075494 & 1.42316670 & MS\\
16.5529 & 12.7633325405 & 1.41866050 & MS\\
20.278 & 15.5223943689 & 1.41364795 & MS\\
24.8414 & 18.9136032758 & 1.40838925 & MS\\
32.562 & 24.6670488835 & 1.40131488 & MS\\
39.8898 & 30.1384093228 & 1.39612751 & MS\\
48.8667 & 36.8490548018 & 1.39113450 & MS\\
\bottomrule\end{tabular}\end{center}
\textbf{Strong-coupling fits} ($\sigma=s^{-1/2}$; coefficients of $1,\sigma,\sigma^2,\dots$ for
free fits, of $\sigma,\sigma^2,\dots$ after subtracting $4/3$ for fixed fits):
\begin{center}\small\begin{tabular}{llll}\toprule fit & $s_{\min}$ & terms & coefficients\\\midrule
free & 15 & 4 & 1.33285, 0.46502, -0.30670, -0.67069\\
free & 15 & 5 & 1.33382, 0.44473, -0.14893, -1.20874, 0.67936\\
$a_1=4/3$ & 15 & 3 & 0.45752, -0.26849, -0.73442\\
$a_1=4/3$ & 15 & 4 & 0.45483, -0.22703, -0.94389, 0.34674\\
free & 20 & 4 & 1.33309, 0.46096, -0.28422, -0.71170\\
free & 20 & 5 & 1.33361, 0.44934, -0.18722, -1.06855, 0.48851\\
$a_1=4/3$ & 20 & 3 & 0.45692, -0.26188, -0.75242\\
$a_1=4/3$ & 20 & 4 & 0.45554, -0.23883, -0.87926, 0.23015\\
free & 25 & 4 & 1.33321, 0.45884, -0.27176, -0.73598\\
free & 25 & 5 & 1.33344, 0.45342, -0.22291, -0.93062, 0.28970\\
$a_1=4/3$ & 25 & 3 & 0.45666, -0.25873, -0.76166\\
$a_1=4/3$ & 25 & 4 & 0.45591, -0.24528, -0.84159, 0.15735\\
\bottomrule\end{tabular}\end{center}
\textbf{Global approximant} (\texttt{strong/global2.py}). With $t=(1+s/\kappa)^{-1/2}$,
\[ m(s)\approx\frac{3s}{4}+\sum_{k=1}^{10}c_k\,t^k , \]
which is analytic at $s=0$ and expands in $s^{-1/2}$ at infinity. Imposed exactly: $m(0)=3/2$,
$m'(0)=1$, $c_1\sqrt\kappa=4/3$, $c_2\kappa=0.456527$; the other six $c_k$ and
$\kappa=0.3562$ are fitted. Maximum relative error 1.1e-05 on
$0<s\le49$. Coefficients: 2.233910, 1.281505, -0.714086, 6.058693, -75.205525, 230.189503, -352.315510, 303.962417, -142.241360, 28.250452.
It is an interpolant, not a closed form; it does not reproduce the complex singularity at
$|s_*|=0.055$, so its Taylor coefficients beyond $s^1$ are not those of $m(s)$.
